\documentclass[UTF8,a4paper,11pt]{ctexart}
\usepackage[margin=22mm,headheight=14pt,footskip=12mm]{geometry}
\usepackage{amsmath,amssymb,array,longtable,booktabs,tikz,needspace,fancyhdr,hyperref}
\usepackage{microtype}
\hypersetup{colorlinks=false,pdfborder={0 0 0},pdfauthor={Good Learning},bookmarks=false}
\IfFontExistsTF{MiSans}{\setCJKmainfont[BoldFont=MiSans Bold]{MiSans}\setCJKsansfont[BoldFont=MiSans Bold]{MiSans}}{}
\newcolumntype{P}[1]{>{\raggedright\arraybackslash}p{#1}}
\setlength{\LTleft}{0pt}\setlength{\LTright}{0pt}
\setlength{\parindent}{0pt}\setlength{\parskip}{.65em}
\linespread{1.18}
\raggedbottom
\setcounter{secnumdepth}{0}
\ctexset{section={format=\Large\bfseries\sffamily,beforeskip=1.35em,afterskip=.65em},
subsection={format=\large\bfseries\sffamily,beforeskip=1em,afterskip=.45em}}
\newcommand{\chaptertitle}[1]{\Needspace{8\baselineskip}\noindent{\LARGE\bfseries\sffamily #1}\par\vspace{4mm}\hrule\vspace{6mm}}
\newcommand{\CompactQuestion}[2]{\par\noindent\textbf{#1.}\ #2\par\vspace{.8em}}
\newcommand{\PracticeQuestion}[3]{\par\Needspace{4\baselineskip}\noindent\begin{minipage}{\linewidth}\textbf{#1.}\ #2\par\noindent\rule{0pt}{#3mm}\par\end{minipage}\par}
\newcommand{\CompactQuestionColon}[2]{\par\noindent\textbf{#1：}\ #2\par\vspace{.8em}}
\newcommand{\PracticeQuestionColon}[3]{\par\Needspace{4\baselineskip}\noindent\begin{minipage}{\linewidth}\textbf{#1：}\ #2\par\noindent\rule{0pt}{#3mm}\par\end{minipage}\par}
\pagestyle{fancy}\fancyhf{}\fancyfoot[C]{\thepage}\renewcommand{\headrulewidth}{0pt}
\begin{document}
\thispagestyle{empty}
\noindent{\sffamily\footnotesize 每日课程}\hfill{\sffamily\footnotesize 学习导图 · 宽松版}\par
\vspace{7mm}
\noindent{\sffamily\fontsize{25}{30}\selectfont\bfseries 第 12 天 · 上采样、检测与分割}\par
\vspace{4mm}\noindent\rule{\linewidth}{.85pt}\par\vspace{6mm}
\begin{center}
\begin{tikzpicture}[x=1cm,y=1cm,>=stealth,every node/.style={font=\sffamily}]
\path[use as bounding box] (-8,-18.6) rectangle (8,.2);
\node[font=\sffamily\Large\bfseries] at (0,-.9) {从特征网格到框和像素标签};
\node[font=\sffamily\small] at (0,-1.68) {左侧看新格怎么填，右侧看两种视觉任务交什么结果};
\node[anchor=west,font=\sffamily\large\bfseries] at (-7.4,-3.35) {网格如何放大};
\draw[line width=.65pt] (-7.4,-3.87) -- ++(6.7,0);
\node[anchor=north west,font=\sffamily\bfseries] at (-6.75,-4.65) {上采样让网格变密};
\node[anchor=north west,align=left,font=\sffamily\small] at (-6.75,-5.430000000000001) {先辨认题目给的是坐标还是核\\[.36em]再选方法与边界规则};
\draw[black!55] (-7.4, -4.9) -- ++(.43,0);
\node[anchor=north west,font=\sffamily\bfseries] at (-6.75,-8.55) {双线性插值};
\node[anchor=north west,align=left,font=\sffamily\small] at (-6.75,-9.33) {在两行分别插值\\[.36em]再按纵向比例加权};
\draw[black!55] (-7.4, -8.8) -- ++(.43,0);
\node[anchor=north west,font=\sffamily\bfseries] at (-6.75,-12.45) {转置卷积};
\node[anchor=north west,align=left,font=\sffamily\small] at (-6.75,-13.229999999999999) {输入乘核，按步长铺放\\[.36em]重叠加总，按填充参数裁边};
\draw[black!55] (-7.4, -12.7) -- ++(.43,0);
\draw[black!55] (-7.4,-4.90) -- (-7.4,-12.70);
\node[anchor=west,font=\sffamily\large\bfseries] at (0.7,-3.35) {视觉任务输出什么};
\draw[line width=.65pt] (0.7,-3.87) -- ++(6.7,0);
\node[anchor=north west,font=\sffamily\bfseries] at (1.35,-4.65) {共享特征图};
\node[anchor=north west,align=left,font=\sffamily\small] at (1.35,-5.430000000000001) {骨干网络提取空间特征\\[.36em]检测从网格位置出发};
\draw[black!55] (0.7, -4.9) -- ++(.43,0);
\node[anchor=north west,font=\sffamily\bfseries] at (1.35,-8.55) {检测给出框和类别};
\node[anchor=north west,align=left,font=\sffamily\small] at (1.35,-9.33) {RPN 提议候选框\\[.36em]第二阶段判类别、修正框};
\draw[black!55] (0.7, -8.8) -- ++(.43,0);
\node[anchor=north west,font=\sffamily\bfseries] at (1.35,-12.45) {分割给出像素标签};
\node[anchor=north west,align=left,font=\sffamily\small] at (1.35,-13.229999999999999) {语义分割区分类别\\[.36em]实例分割还区分同类个体};
\draw[black!55] (0.7, -12.7) -- ++(.43,0);
\draw[black!55] (0.7,-4.90) -- (0.7,-12.70);
\node[font=\sffamily\small] at (0,-16.35) {10 分钟热身 · 25 分钟讲解 · 40 分钟练习 · 15 分钟订正};
\end{tikzpicture}
\end{center}
\clearpage

\chaptertitle{第 12 天 · 上采样、检测与分割}
\textbf{2026 年 11 月 26 日｜K09、K13–K15｜90 分钟}

把一张小图放大时，新格子该填什么数？看到一张有两只猫的图时，模型该画两个框，还是把每只猫的轮廓涂出来？今天先学两种把网格变大的方法，再看检测和分割各交出什么结果。

\section{90 分钟安排}
\begin{itemize}
\item 10 分钟 · H1 用 5 分钟，H2 和 H3 用余下 5 分钟口答。
\item 25 分钟 · 读三个例题，做 A1、B1、C1。用到的前置是矩阵行列、加权平均和卷积窗口。
\item 40 分钟 · A2、B2、C2 各 8 分钟，M1 用 12 分钟，M2 用 4 分钟。
\item 15 分钟 · 对照解析，合上答案重做最先出错的一步，再口述检测和两类分割的区别。
\end{itemize}
H1–H3 来自本课程错题本第 23–25 条的复做题。其余均为新编题，覆盖备考计划中计算 9、22–24 所对应的考点，并穿插 K09 的回归树复习。

\subsection{举一反三 · 热身}
\PracticeQuestion{H1}{输入为 $X=\begin{pmatrix}2&1&3\\4&0&2\\1&5&6\end{pmatrix}$，核为 $K=\begin{pmatrix}1&0\\0&-1\end{pmatrix}$，步长 2，四边各补 1 圈零，偏置 0。先画出补零后的矩阵，再算四格输出，圈出碰到补零圈的窗口。坐标按补零后的矩阵从 $(1,1)$ 起算。}{45}
\PracticeQuestion{H2}{单通道 $5\times5$ 图用一个带偏置的 $3\times3$ 核，步长 1、无填充。一个输出读几个像素？相邻两个窗口共用几项参数？各自对应哪个概念？}{12}
\PracticeQuestion{H3}{图像的位置网格从 $100\times100$ 变为 $50\times50$，亮度级数不变，改变的是哪一步？若位置不动，亮度级从 256 改为 16 呢？}{12}
\clearpage
\section{单元 1 · 双线性插值，用四个邻点填新格子}
上采样会增加空间位置。双线性插值给每个新位置算一个加权平均。先由题设的坐标规则找到它在原网格上的位置。若题目给了原图坐标，就直接用该坐标。格子外的位置要依题意处理边界。

四角写作 $\begin{pmatrix}a&b\\c&d\end{pmatrix}$。位置在格子内，从左到右走了 $u$，从上到下走了 $v$，其中 $0\le u,v\le1$。先在两行分别算插值，再在两行之间插一次：

\[t=(1-u)a+ub,\qquad q=(1-u)c+ud.\]
\[z=(1-v)t+vq.\]
左右用的是 $u$，上下用的是 $v$。四角最终权重都非负且和为 1，结果应落在四个数的最小值和最大值之间。

\subsection{完整例题 · 先沿横向走，再沿纵向走}
四角是 $\begin{pmatrix}2&6\\8&14\end{pmatrix}$，求 $u=\frac14,\ v=\frac12$ 处的值。

\begin{enumerate}
\item 上边向右只走四分之一，左边的 2 占四分之三，得 $t=\frac34(2)+\frac14(6)=3$。
\item 下边用同样的横向比例，得 $q=\frac34(8)+\frac14(14)=9.5$。
\item 纵向走了一半，取两值的平均，得 $z=(3+9.5)/2=6.25$。它落在 2 和 14 之间。若 $v=0$，该点在上边，结果只由上边两角决定。
\end{enumerate}
\subsection{举一反三}
\PracticeQuestion{A1}{沿用例题的四角和 $u=\frac14$。把 $v$ 改成 $\frac14$，已知 $t=3,\ q=9.5$。补最后一步的算式。}{23}
\PracticeQuestion{A2}{四角改为 $\begin{pmatrix}0&8\\4&20\end{pmatrix}$，求 $u=\frac34,\ v=\frac14$ 处的值。再令 $v=1$，保持 $u$ 不变，求新值并说明哪些角的权重变为 0。指出与例题相比哪一步的权重变了。}{26}
\clearpage
\section{单元 2 · 转置卷积，每个输入铺一块，重叠就相加}
普通卷积从一个窗口取数，汇成一个输出。转置卷积可从输入出发想：每格输入乘上同一个核，把这块贡献铺到输出上。同一格收到多块，就相加。核的权重通常由训练学得。它对应卷积线性运算的矩阵转置，通常无法还原已经丢掉的原图信息。

本课使用与昨天的互相关相配的转置约定，按给定核方向铺放，偏置 0、膨胀率 1、额外输出填充 0。输入每边长 $N$、步长 $S$、核边长 $K$ 时，先从第 1 格开始铺，下一块起点相隔 $S$ 格，完整范围长 $(N-1)S+K$。若题设转置卷积的填充参数为 $P$，按这个算法在完整范围两端各裁 $P$ 格，保留长度：

\[O=(N-1)S-2P+K.\]
对高和宽各算一次。$S<K$ 时相邻的块在该方向重叠。题目若改了膨胀率或额外输出填充，尺寸公式要相应改变。

\subsection{完整例题 · 把两个输入铺到输出上}
输入只有一行 $(2,3)$，核只有一行 $(1,2,1)$。横向步长 2，高度方向维持 1，$P=0$。求输出。

\begin{enumerate}
\item 输出宽 $(2-1)\cdot2+3=5$，因为最后一块从第 3 格起，还要占 3 格。
\item 输入 2 的贡献是 $(2,4,2)$，放第 1–3 格。输入 3 的贡献是 $(3,6,3)$，放第 3–5 格。
\item 第 3 格收到 $2+3=5$，其余各读一块。输出为 $(2,4,5,6,3)$，共 $1\times5$。检查总和 $20=(2+3)(1+2+1)$。这项检查适用于本题的无裁剪、零偏置条件。
\end{enumerate}
\subsection{举一反三}
\PracticeQuestion{B1}{输入改成 $(1,4)$，沿用例题的核和设置。已知两块贡献分别为 $(1,2,1)$ 和 $(4,8,4)$。写出重叠格的位置和它的值。}{24}
\PracticeQuestion{B2}{输入 $X=\begin{pmatrix}1&2\\0&3\end{pmatrix}$，核 $K=\begin{pmatrix}1&2\\-1&1\end{pmatrix}$，高宽步长都是 1，$P=0$。标出四块的左上角，算出完整输出及中央一格的四项贡献。与例题相比，说出重叠有什么变化。若高宽的 $P$ 都改成 1，保留几行几列？给出留下的值。}{34}
\clearpage
\section{单元 3 · 从整张图的标签，到框，再到逐像素标签}
假设照片里有两只猫。分类可以说「图里有猫」。检测为两只猫分别给出框和类别。语义分割把两只猫的像素都标作「猫」。实例分割还会区分这些像素属于猫 1 还是猫 2。

Faster R-CNN 是两阶段检测器。骨干网络先生成共享特征图。区域建议网络（RPN）在每个特征图位置放若干预设参考框（Anchor），分别给出前景评分并预测框的四个偏移量（中心横纵位置与宽高的调整）。将偏移量应用到 Anchor，再筛选，得到候选框。第二阶段从共享特征图提取每个候选区域的特征，判断具体类别并进一步修正框。RPN 的前景分数回答「这里像不像有物体」。

\subsection{完整例题 · RPN 有几个框、几个数}
特征图高 2、宽 3，每处有 3 个 Anchor。本题每个 Anchor 用 2 个分类分数（前景、背景）和 4 个回归量。忽略筛选，按高、宽、通道的顺序写输出规模。

\begin{enumerate}
\item 一共有 $2\cdot3=6$ 个位置，每处 3 个框，原始 Anchor 数是 $6\cdot3=18$。
\item 每处有 $3\cdot2=6$ 个分类分数，形状是 $2\times3\times6$，全图 36 个分数。
\item 每处有 $3\cdot4=12$ 个回归量，形状是 $2\times3\times12$，全图 72 个量。交叉检查：$18\cdot4=72$。如果题设采用单前景分数，分类通道是 $3\cdot1=3$，框数和四个偏移量的计数仍按上述规则。
\end{enumerate}
\subsection{举一反三}
\PracticeQuestion{C1}{例题改为每处 4 个 Anchor，每框仍给两个分类分数。已知全图有 24 个 Anchor，补分类分支的通道数和全图分数数目。}{27}
\PracticeQuestion{C2}{特征图改为 $3\times4$，每处 2 个 Anchor，每框给 1 个前景分数和 4 个回归量。求原始框数、两分支各自的形状与元素数。指出哪条计数规则因分数约定而变。若筛选后只保留 5 个候选框，用一句话说明第二阶段接下来做什么。}{35}
\clearpage
\section{独立练习}
\textbf{先写出过程，再翻到后面的解析。}

\PracticeQuestion{M1}{一张路面图里有两辆紧挨的汽车。模型有一张 $2\times3$ 的低分辨率检测特征图，每处 2 个 Anchor，RPN 每框输出两个分类分数和四个偏移量。另一条分割分支在一个局部方格有四个分数 $\begin{pmatrix}0&2\\4&10\end{pmatrix}$，需要估计 $u=\frac12,\ v=\frac34$ 处的分数。（1） 求 RPN 的原始框数及两个分支的全图元素数。（2） 算新位置的分数，写出两个方向的权重。（3） 如果要分别数清两辆车各占哪些像素，应选哪种任务？只看检测框或语义分割分别缺哪项信息？}{40}
\PracticeQuestion{M2}{回顾 K09。回归树某叶包含三个样本，目标值是 $2,4,9$，采用平方误差。该叶输出多少？同一组 $2,4,9$ 若作为一行特征做最大池化，且窗口覆盖整行，结果是多少？各用一句话说明为什么。}{34}
\clearpage
\section{答案与解析 · 热身和单元题}
\textbf{H1.} 补零后的五行为 $(0,0,0,0,0)$、$(0,2,1,3,0)$、$(0,4,0,2,0)$、$(0,1,5,6,0)$、$(0,0,0,0,0)$。每边有 $\lfloor(5-2)/2\rfloor+1=2$ 格。左上角按行依次为 $(1,1),(1,3),(3,1),(3,3)$。每窗取左上减右下，依次为 $0-2=-2$、$0-3=-3$、$0-1=-1$、$0-6=-6$。输出 $\begin{pmatrix}-2&-3\\-1&-6\end{pmatrix}$，前三个窗口碰到补零圈。右下窗的左上是原图里的 0。

\textbf{H2.} 一个输出只读窗口内 9 个像素，这是局部连接。各窗口使用同一组 9 个核权重和 1 个偏置，共 10 个参数，这是权值共享。

\textbf{H3.} 改变读数位置的网格是改变采样。位置固定、亮度级数变化是改变量化。

\textbf{A1.} 横向位置相同，因此 $t=3,\ q=9.5$。最后用纵向权重，$z=\frac34(3)+\frac14(9.5)=4.625$。

\textbf{A2.} 右角占 $\frac34$，得 $t=\frac14(0)+\frac34(8)=6$、$q=\frac14(4)+\frac34(20)=16$。下边占 $\frac14$，得 $z=\frac34(6)+\frac14(16)=8.5$。与例题相比，横向和纵向比例都变了。$v=1$ 时，上边两角的权重都为 0，直接得 $q=16$。

\textbf{B1.} 第二块从第 3 格起，第 3 格叠加两块，得 $1+4=5$。整行是 $(1,2,5,8,4)$。

\textbf{B2.} 每边 $(2-1)\cdot1+2=3$。四块按行从 $(1,1),(1,2),(2,1),(2,2)$ 铺起，分别为 $K,2K,0K,3K$。中央 $(2,2)$ 依次收到 $1,-2,0,3$，相加得 2。整个输出为 $\begin{pmatrix}1&4&4\\-1&2&8\\0&-3&3\end{pmatrix}$。例题只在横向叠两块，这里中心可以叠四块。无裁剪时总和为 $18=(1+2+0+3)(1+2-1+1)$。$P=1$ 时四边各裁 1 格，只保留 $1\times1$ 的中间值 2。

\textbf{C1.} 每处 4 框，各给两个分数，因此 8 个分类通道。$2\cdot3\cdot8=48$，也等于 $24\cdot2$。

\textbf{C2.} 12 个位置，每处 2 框，共 24 框。分类每处 $2\cdot1=2$ 个，形状 $3\times4\times2$，共 24 个。回归每处 $2\cdot4=8$ 个，形状 $3\times4\times8$，共 96 个。单分数约定影响分类通道数。第二阶段根据 5 个候选区域的特征判断类别并微调框。

\clearpage
\subsection{混合题答案}
\textbf{M1.} （1） $2\cdot3\cdot2=12$ 框。分类共 $12\cdot2=24$ 个分数，回归共 $12\cdot4=48$ 个量。（2） 横向权重各 $\frac12$，两行分别插值得 1 和 7。纵向权重上 $\frac14$、下 $\frac34$，新分数为 $\frac14(1)+\frac34(7)=5.5$，介于 0 和 10 之间。（3） 用实例分割。检测框缺少逐像素轮廓，语义分割有汽车像素标签，但缺少每辆车各自的身份标签。

\textbf{M2.} 回归叶取均值 $(2+4+9)/3=5$，因为它使叶内平方误差和最小。最大池化按定义取窗口内最大值 9。

\section{完成标准}
\begin{itemize}
\item 独立手算出一个双线性插值值，并写清横向、纵向权重。能在转置卷积题中圈出重叠处，写出参与相加的各项。
\item 用两辆车或两只猫举例，说清检测、语义分割和实例分割的输出区别。能按题设的每框分数数目，独立写出 RPN 两个分支的输出数量。
\item 发来 H1、A2、B2、C2 或 M1 的步骤和用时即可批改。卡住时可只发题号和已经写到的那一步。
\end{itemize}
转置卷积尺寸选读：\href{https://arxiv.org/abs/1603.07285}{Dumoulin 与 Visin}。Faster R-CNN 及双分数 RPN 选读：\href{https://proceedings.neurips.cc/paper/2015/file/14bfa6bb14875e45bba028a21ed38046-Paper.pdf}{Ren 等，原论文图 1}。

\end{document}
